\section*{Abstract}
\noindent\textbf{Research Question/Issue:} In a Korean captive spin-off a listed parent places a business in a wholly owned subsidiary and its controlling owner later decides, without a shareholder vote, whether to keep, sell or list it; minority shareholders lose when it is disposed of. Was the intention to dispose legible in the division filing before the loss, did the market price it, and what did the 2022 disclosure reform change?
\par\noindent\textbf{Research Findings/Insights:} Two large language models read \NScored{} captive spin-off filings of 2001--2026 as of their dates. Their forecast separates subsidiaries later disposed of from those retained (AUC \AUCPooledLLM{} against \AUCPooledFin{} for the filing's financial characteristics) and survives the masking of all names and dates. The announcement return rises with the forecast while the market-adjusted return over the following year falls with it (Spearman \SpearmanScoreBhar{}); the loss is realised in parents that dispose within a year and in parents with small controlling blocks, and the forecast falls with the block. After the reform, protection statements rose from \ProtectPre{} to \ProtectPost{} percent of filings and the announcement return from \CarPre{} to \CarPost{} percent; the forecast and the early-disposal rate did not change.
\par\noindent\textbf{Theoretical/Academic Implications:} The principal--principal literature infers the controller's intention from the expropriating outcome. Here it was observable in the controller's own disclosure beforehand, and minority investors priced it in the wrong direction.
\par\noindent\textbf{Practitioner/Policy Implications:} The information minority shareholders needed was disclosed and unread; mandating disclosure changed the filings and their reception first. Protection should attach to the disposal.
\par\medskip\noindent\textbf{Keywords:} captive spin-offs; controlling shareholders; minority shareholder protection; corporate disclosure; business groups; Korea